%% toc: P3 25 25.4
\section{Proofs for saturated elementary extensions}
\label{sec:annex-p03-w03-p03-010-saturated-extension}

Throughout, the Axiom of Choice holds in the ambient set theory and
cardinals are initial ordinals.  For a structure \(\mathcal M\) write
\(M=\dom\mathcal M\).  For \(A\subseteq M\), every \(\Lang(A)\)-formula
\(\varphi(\bar x)\) has the form
\(\chi(\bar x,c_{a_1},\ldots,c_{a_m})\) with \(\chi(\bar x,\bar y)\) an
\(\Lang\)-formula and \(a_1,\ldots,a_m\in A\); then
\(\mathcal M_A\entails\varphi(\bar b)\) means
\(\mathcal M\entails\chi(\bar b,a_1,\ldots,a_m)\).  This reading is used
without further mention.  For \(A\subseteq M\) the constants \(c_a\)
(\(a\in A\)) of \(\Lang(A)\) are the same symbols as in \(\Lang(M)\).

\begin{annexitem}{annex:p03-w03-p03-010-type-facts}{thm:p03-w03-p03-010-type-facts}{Consistency, completion and transfer of types}
\begin{proof}
Let \(\bar c=(c_1,\ldots,c_n)\) be new constants and put
\[
  \Gamma=\operatorname{Th}_{\Lang(A)}(\mathcal M_A)\cup\{\varphi(\bar c):\varphi\in p\}.
\]
By definition, \(p\) is a type over \(A\) consistent with \(\mathcal M\)
exactly when \(\Gamma\) is satisfiable.

(1) Suppose \(\Gamma\) has a model \(\mathcal K\), and let
\(\varphi_1,\ldots,\varphi_k\in p\).  If \(k=0\), any \(n\)-tuple of the
nonempty domain \(M\) realizes the empty subset.  If \(k\geq1\), let
\(\sigma\) be the \(\Lang(A)\)-sentence
\(\exists\bar x(\varphi_1\land\cdots\land\varphi_k)\).  The tuple
\(\bar c^{\mathcal K}\) witnesses \(\mathcal K\entails\sigma\).  If
\(\mathcal M_A\) did not satisfy \(\sigma\), then \(\neg\sigma\) would be
in \(\operatorname{Th}_{\Lang(A)}(\mathcal M_A)\subseteq\Gamma\), and \(\mathcal K\)
would satisfy both \(\sigma\) and \(\neg\sigma\).  So
\(\mathcal M_A\entails\sigma\), and a witnessing tuple \(\bar b\) from
\(M\) realizes \(\{\varphi_1,\ldots,\varphi_k\}\) in \(\mathcal M\).

Conversely, suppose every finite subset of \(p\) is realized in
\(\mathcal M\).  A finite \(\Gamma_0\subseteq\Gamma\) contains only
finitely many sentences \(\varphi_1(\bar c),\ldots,\varphi_k(\bar c)\)
with \(\varphi_i\in p\).  Choose \(\bar b\) from \(M\) realizing
\(\{\varphi_1,\ldots,\varphi_k\}\), and expand \(\mathcal M_A\) by
interpreting \(\bar c\) as \(\bar b\).  This expansion satisfies
\(\operatorname{Th}_{\Lang(A)}(\mathcal M_A)\) and each \(\varphi_i(\bar c)\), hence
\(\Gamma_0\).  By first-order compactness (\cref{thm:fol-compactness},
clause~1, in the language \(\Lang(A)\cup\{c_1,\ldots,c_n\}\), which may
be uncountable; the Axiom of Choice is assumed), \(\Gamma\) is satisfiable.

(2) Let \(\mathcal K\entails\Gamma\), and let \(q\) be the set of all
\(\Lang(A)\)-formulas \(\varphi(\bar x)\), free variables among \(\bar x\),
with \(\mathcal K\entails\varphi(\bar c)\).  Then \(p\subseteq q\).  For
each such \(\varphi\), exactly one of \(\varphi(\bar c)\) and
\(\neg\varphi(\bar c)\) is true in \(\mathcal K\), so exactly one of
\(\varphi,\neg\varphi\) is in \(q\).  The structure \(\mathcal K\) is a model
of \(\operatorname{Th}_{\Lang(A)}(\mathcal M_A)\cup\{\varphi(\bar c):\varphi\in q\}\),
so \(q\) is a complete type over \(A\) consistent with \(\mathcal M\).

(3) Let \(\sigma=\chi(c_{a_1},\ldots,c_{a_m})\) be an
\(\Lang(A)\)-sentence with \(\chi\) an \(\Lang\)-formula.  Since
\(a_1,\ldots,a_m\in M\) and \(\mathcal M\preccurlyeq\mathcal M'\),
\begin{align*}
  \mathcal M_A\entails\sigma
  &\iff\mathcal M\entails\chi(a_1,\ldots,a_m)\\
  &\iff\mathcal M'\entails\chi(a_1,\ldots,a_m)
  \iff\mathcal M'_A\entails\sigma .
\end{align*}
So the two theories are equal.  Consistency of \(p\) with a structure
depends only on this theory, and completeness of \(p\) does not depend on
the structure at all.  For realization, let \(\bar b\) be from \(M\) and
\(\varphi=\chi(\bar x,c_{a_1},\ldots,c_{a_m})\in p\); then
\(\mathcal M\entails\chi(\bar b,\bar a)\) iff
\(\mathcal M'\entails\chi(\bar b,\bar a)\), again by elementarity.
\end{proof}
\end{annexitem}

\begin{annexitem}{annex:p03-w03-p03-010-copy}{thm:p03-w03-p03-010-copy}{Replacing an elementary embedding by an elementary extension}
\begin{proof}
Write \(S=\bigcup\bigcup M\) and \(s=s_{\mathcal M}=\{u\in S:u\notin u\}\).
If \(s\in S\), then \(s\in s\) iff \(s\notin s\); hence \(s\notin S\).
Put \(X=\{(y,s):y\in N^*\setminus h[M]\}\), where \(N^*=\dom\mathcal
N^*\).  If \((y,s)=\{\{y\},\{y,s\}\}\) belonged to \(M\), then
\(\{y,s\}\in\bigcup M\) and \(s\in S\), which is impossible.  So
\(M\cap X=\varnothing\).

The map \(h\) is injective, since \(a\ne b\) in \(\mathcal M\) gives
\(h(a)\ne h(b)\) in \(\mathcal N^*\) by elementarity.  Define
\(F:N^*\to M\cup X\) by \(F(h(a))=a\) for \(a\in M\) and \(F(y)=(y,s)\)
for \(y\notin h[M]\).  Each clause is injective, and their ranges \(M\) and
\(X\) are disjoint, so \(F\) is a bijection onto \(M\cup X\).  Make
\(D=M\cup X\) into an \(\Lang\)-structure \(\mathcal N\) by transport:
\(c^{\mathcal N}=F(c^{\mathcal N^*})\) for constants,
\(f^{\mathcal N}(\bar d)=F(f^{\mathcal N^*}(F^{-1}\bar d))\) for function
symbols, and \(\bar d\in R^{\mathcal N}\) iff
\(F^{-1}\bar d\in R^{\mathcal N^*}\) for relation symbols, with \(F^{-1}\)
applied coordinatewise.

By induction on terms, \(F(t^{\mathcal N^*}[\bar e])=t^{\mathcal
N}[F\bar e]\).  By induction on formulas, \(\mathcal
N^*\entails\varphi(\bar e)\) iff \(\mathcal N\entails\varphi(F\bar e)\):
equations hold because \(F\) is injective, relation instances hold by the
transported definition, the Boolean steps are immediate, and in the
existential step a witness \(d\) for \(\mathcal N^*\) gives the witness
\(F(d)\) for \(\mathcal N\), while every \(d'\) in \(D\) has the form
\(F(d)\) because \(F\) is onto.

Now let \(\varphi\) be an \(\Lang\)-formula and \(\bar a\) a tuple from
\(M\).  By elementarity of \(h\) and the preceding paragraph,
\[
  \mathcal M\entails\varphi(\bar a)
  \iff\mathcal N^*\entails\varphi(h\bar a)
  \iff\mathcal N\entails\varphi(F h\bar a)
  =\varphi(\bar a).
\]
Applied to the atomic formulas \(c=y\), \(f(\bar y)=z\) and \(R(\bar y)\),
this shows that the symbols of \(\Lang\) are interpreted in \(\mathcal M\) as
the restrictions to \(M\) of their interpretations in \(\mathcal N\), so
\(\mathcal M\) is a substructure of \(\mathcal N\); applied to all formulas,
it shows \(\mathcal M\preccurlyeq\mathcal N\).  Finally \(F(h(a))=a\) by
definition.
\end{proof}
\end{annexitem}

\begin{annexitem}{annex:p03-w03-p03-010-one-step}{thm:p03-w03-p03-010-one-step}{Realizing every type over every subset in one elementary extension}
\begin{proof}
\emph{The task set.}  Let \(P=P_{\mathcal M}\) be the set of triples
\(\pi=(A,n,p)\) with \(A\subseteq M\), \(n<\omega\), and \(p\) a complete
type over \(A\) in the variables \(x_1,\ldots,x_n\) consistent with
\(\mathcal M\).  It is a subset of
\(\mathcal P(M)\times\omega\times\mathcal P(\operatorname{Form}(\Lang(M)))\)
and hence a set.  (Writing the arity into the triple only makes the pairs
\((A,p)\) of the statement unambiguous.)  Let
\(\mu=\max(\abs\Lang,\abs M,\abs P,\aleph_0)\).

\emph{The theory.}  For each \(\pi=(A,n,p)\in P\) take new constants
\(d_{\pi,1},\ldots,d_{\pi,n}\), distinct for distinct pairs \((\pi,i)\) and
not in \(\Lang(M)\).  Let \(\Lang'\) be \(\Lang(M)\) with all of them.  For
\(\varphi(\bar x)\in p\), let \(\varphi(\bar d_\pi)\) be the
\(\Lang'\)-sentence obtained by substituting \(d_{\pi,i}\) for \(x_i\).
Put
\[
  \Sigma=\operatorname{ElDiag}(\mathcal M)\cup
  \{\varphi(\bar d_\pi):\pi=(A,n,p)\in P,\ \varphi\in p\}.
\]
By \cref{thm:p06-cardinal-arithmetic-laws},
\(\abs{\Lang'}\leq\abs\Lang+\abs M+\abs P\cdot\aleph_0\leq\mu\).

\emph{Finite satisfiability.}  A finite \(\Sigma_0\subseteq\Sigma\)
mentions constants \(\bar d_\pi\) for only finitely many
\(\pi_1,\ldots,\pi_k\in P\), and for each \(j\) only a finite set
\(F_j\subseteq p_{\pi_j}\) of formulas.  By
\cref{thm:p03-w03-p03-010-type-facts}~(1), some tuple \(\bar b_j\) from
\(M\) realizes \(F_j\) in \(\mathcal M\).  Expand \(\mathcal M_M\) by
interpreting \(\bar d_{\pi_j}\) as \(\bar b_j\) for \(j\leq k\) and every
other new constant as one fixed element of \(M\).  This expansion satisfies
\(\operatorname{ElDiag}(\mathcal M)\), because it expands \(\mathcal M_M\),
and satisfies each \(\varphi(\bar d_{\pi_j})\) with \(\varphi\in F_j\).  So
it is a model of \(\Sigma_0\).  By \cref{thm:fol-compactness}, clause~1,
\(\Sigma\) has a model \(\mathcal K\).

\emph{Bounding the size.}  By the downward L\"owenheim--Skolem theorem
(\cref{thm:downward-ls}, with the empty set of named elements) there is
\(\mathcal K'\preccurlyeq\mathcal K\) with
\(\abs{\dom\mathcal K'}\leq\max(\abs{\Lang'},\aleph_0)\leq\mu\).  It
satisfies every sentence of \(\Sigma\).  Choose an injection
\(\iota:\dom\mathcal K'\to\mu\) and transport \(\mathcal K'\) along
\(\iota\), as in the proof of \cref{thm:p03-w03-p03-010-copy}; the result
satisfies the same sentences.  So we may assume
\(\dom\mathcal K'\subseteq\mu\).

\emph{The embedding.}  Let \(\mathcal N^*\) be the \(\Lang\)-reduct of
\(\mathcal K'\) and put \(h(a)=c_a^{\mathcal K'}\) for \(a\in M\).  For an
\(\Lang\)-formula \(\chi\) and a tuple \(\bar a\) from \(M\), if
\(\mathcal M\entails\chi(\bar a)\) then
\(\chi(c_{a_1},\ldots,c_{a_m})\in\operatorname{ElDiag}(\mathcal M)\), so
\(\mathcal N^*\entails\chi(h\bar a)\); if not, then
\(\neg\chi(c_{a_1},\ldots,c_{a_m})\) is in the diagram and
\(\mathcal N^*\) fails \(\chi(h\bar a)\).  So \(h\) is an elementary
embedding.  (This is the converse half of \cref{thm:diagram-tests},
clause~2.)  By \cref{thm:p03-w03-p03-010-copy} there are
\(\mathcal M'\succcurlyeq\mathcal M\) and an isomorphism
\(F:\mathcal N^*\to\mathcal M'\) with \(F\circ h\) the identity on \(M\) and
\[
  \dom\mathcal M'=M\cup\{(y,s_{\mathcal M}):y\in\dom\mathcal
  N^*\setminus h[M]\}\subseteq M\cup(\mu\times\{s_{\mathcal M}\}).
\]

\emph{Realization.}  Let \(\pi=(A,n,p)\in P\) and put
\(e_i=F(d_{\pi,i}^{\mathcal K'})\).  For
\(\varphi=\chi(\bar x,c_{a_1},\ldots,c_{a_m})\in p\), the sentence
\(\varphi(\bar d_\pi)\) is in \(\Sigma\), so
\[
  \mathcal N^*\entails\chi(\bar d_\pi^{\mathcal K'},h(a_1),\ldots,h(a_m)),
\]
because \(c_a^{\mathcal K'}=h(a)\).  Applying the isomorphism \(F\) and
\(F(h(a))=a\) gives \(\mathcal M'\entails\chi(\bar e,a_1,\ldots,a_m)\).
So \(\bar e\) realizes \(p\) in \(\mathcal M'\).

\emph{The candidate set.}  All \(\Lang\)-structures with domain a subset of
\(D_{\mathcal M}=M\cup(\mu\times\{s_{\mathcal M}\})\) form a set, since
each is given by its domain and interpretations, which are subsets of
powers of \(D_{\mathcal M}\).  Those that are elementary extensions of
\(\mathcal M\) and realize every \(p\) as above form a subset of it, defined
from \(\mathcal M\) alone, and \(\mathcal M'\) shows that it is nonempty.
\end{proof}
\end{annexitem}

\begin{annexitem}{annex:p03-w03-p03-010-bookkeeping}{thm:p03-w03-p03-010-bookkeeping}{A small set lies in an earlier stage of a chain of length \(\kappa^+\)}
\begin{proof}
\emph{Stages below \(\kappa^+\) are small.}  Let \(\alpha<\kappa^+\).  The
cardinal \(\abs\alpha\) satisfies \(\abs\alpha\leq\alpha<\kappa^+\), and
\(\kappa^+\) is the least cardinal greater than \(\kappa\); cardinals are
linearly ordered under the Axiom of Choice, so \(\abs\alpha\leq\kappa\).

\emph{\(\kappa^+\) is a limit ordinal.}  It is infinite, being greater than
\(\kappa\geq\aleph_0\).  If \(\kappa^+=\beta+1\), then \(\beta\) is
infinite and \(\abs{\beta+1}=\abs\beta\) (send \(\beta\) to \(0\), each
natural number \(n\) to \(n+1\), and fix every other element).  Then
\(\abs{\kappa^+}=\abs\beta\leq\kappa\) by the previous paragraph, which
contradicts \(\kappa<\kappa^+\).

\emph{The bound.}  If \(Y=\varnothing\), take \(\gamma=0\).  Otherwise,
for each \(y\in Y\) choose \(\alpha_y<\kappa^+\) with
\(y\in X_{\alpha_y}\), and let
\(\gamma=\bigcup_{y\in Y}\alpha_y\), the supremum of these ordinals.  It is
an ordinal with \(\alpha_y\leq\gamma\) for every \(y\), and
\(\gamma\subseteq\kappa^+\), so \(\gamma\leq\kappa^+\).  Fix a well-ordering
of \(Y\) and, for each \(y\), an injection \(g_y:\alpha_y\to\kappa\), which
exists because \(\abs{\alpha_y}\leq\kappa\).  For \(\xi\in\gamma\), let
\(y(\xi)\) be the least \(y\) in that well-ordering with
\(\xi\in\alpha_y\), and map \(\xi\mapsto(y(\xi),g_{y(\xi)}(\xi))\).  This
map \(\gamma\to Y\times\kappa\) is injective: equal images have the same
first coordinate \(y\), and \(g_y\) is injective.  Hence
\[
  \abs\gamma\leq\abs Y\cdot\kappa\leq\kappa\cdot\kappa=\kappa
\]
by \cref{thm:p06-cardinal-arithmetic-laws}.  As
\(\abs{\kappa^+}>\kappa\), \(\gamma\ne\kappa^+\), so \(\gamma<\kappa^+\).
Since the sets increase, \(y\in X_{\alpha_y}\subseteq X_\gamma\) for
every \(y\in Y\), that is, \(Y\subseteq X_\gamma\).
\end{proof}
\end{annexitem}

\begin{annexitem}{annex:p03-w03-p03-010-saturated-extension}{thm:p03-w03-p03-010-saturated-extension}{Every structure has a \(\kappa\)-saturated elementary extension}
\begin{proof}
Put \(\lambda=\kappa^+\).  For a structure \(\mathcal K\), let
\(W(\mathcal K)\) be the nonempty set of extensions given by
\cref{thm:p03-w03-p03-010-one-step}, consisting of the elementary extensions
of \(\mathcal K\) with domain inside \(D_{\mathcal K}\) that realize every
complete type in finitely many variables over every subset of
\(\dom\mathcal K\) consistent with \(\mathcal K\).

\emph{Choosing the stages within a set.}  Define sets \(\mathrm{Seq}_\beta\)
for \(\beta\leq\lambda\) by recursion, with no choices.  Let
\(\mathrm{Seq}_0=\{\{(0,\mathcal M)\}\}\).  Let
\(\mathrm{Seq}_{\beta+1}\) be the set of all
\(t\cup\{(\beta+1,\mathcal K)\}\) with \(t\in\mathrm{Seq}_\beta\) and
\(\mathcal K\in W(t(\beta))\); it is a set by Replacement and Union.  For a
limit \(\delta\), let \(\mathrm{Seq}_\delta\) be the set of functions
\(t\) with domain \(\delta+1\) such that
\(t\restriction(\alpha+1)\in\mathrm{Seq}_\alpha\) for every \(\alpha<\delta\),
the structures \(t(\alpha)\) (\(\alpha<\delta\)) are substructures of each
other in increasing order, and \(t(\delta)\) is their union.  For such a
\(t\), the restriction \(t\restriction\delta\) is the union of the
functions \(t\restriction(\alpha+1)\in\mathrm{Seq}_\alpha\), so it is a
subset of the set \(Q_\delta=\bigcup\bigl(\bigcup_{\alpha<\delta}
\mathrm{Seq}_\alpha\bigr)\), and \(t\) is determined by
\(t\restriction\delta\) because \(t(\delta)\) is the union of its values.
Hence \(\mathrm{Seq}_\delta\) is the image of a subset of the power set of
\(Q_\delta\) under the operation
\(r\mapsto r\cup\{(\delta,\mathcal U_r)\}\), where \(\mathcal U_r\) is the
union structure of the increasing structures \(r(\alpha)\),
\(\alpha<\delta\); it is a set by Power Set, Separation and Replacement.  Let
\(Z=\bigcup_{\beta\leq\lambda}\mathrm{Seq}_\beta\), and fix a well-ordering
\(\lhd\) of \(Z\).

Now define \(s_\beta\in\mathrm{Seq}_\beta\) for \(\beta\leq\lambda\) by
recursion, proving at the same time the claim
\[
  (C_\beta)\qquad s_\beta(\alpha)\preccurlyeq s_\beta(\alpha')
  \quad\text{whenever }\alpha\leq\alpha'\leq\beta .
\]
Put \(s_0=\{(0,\mathcal M)\}\); \((C_0)\) is trivial.  Given \(s_\beta\),
let \(s_{\beta+1}\) be the \(\lhd\)-least \(t\in\mathrm{Seq}_{\beta+1}\)
with \(t\restriction(\beta+1)=s_\beta\); such a \(t\) exists because
\(W(s_\beta(\beta))\ne\varnothing\).  Then
\(s_{\beta+1}(\beta)\preccurlyeq s_{\beta+1}(\beta+1)\), and with
\((C_\beta)\) and transitivity of \(\preccurlyeq\) this gives
\((C_{\beta+1})\).  (If \(\mathcal A\preccurlyeq\mathcal B\preccurlyeq
\mathcal C\), a formula with parameters from \(\mathcal A\) has the same
truth value in all three.)  At a limit \(\delta\), the earlier
\(s_\alpha\) extend one another, since each successor step extends the
previous sequence and each limit step extends all earlier ones.  Let
\(u=\bigcup_{\alpha<\delta}s_\alpha\), a function on \(\delta\).  By
\((C_\alpha)\) for \(\alpha<\delta\),
\(\langle u(\alpha):\alpha<\delta\rangle\) is an elementary chain, so its
union \(\mathcal U\) is defined, and \(s_\delta=u\cup\{(\delta,\mathcal
U)\}\) lies in \(\mathrm{Seq}_\delta\).  The elementary chain theorem
(\cref{thm:elementary-chain}) gives \(u(\alpha)\preccurlyeq\mathcal U\) for
\(\alpha<\delta\), which with the earlier claims is \((C_\delta)\).

Write \(\mathcal N_\alpha=s_\lambda(\alpha)\) and
\(\mathcal N=\mathcal N_\lambda\).  This is the construction of
\cref{def:p03-w03-p03-010-construction}: \(\mathcal N_0=\mathcal M\),
\(\mathcal N_{\alpha+1}\in W(\mathcal N_\alpha)\), and unions are taken at
limits.  Since \(\lambda\) is a limit ordinal
(\cref{thm:p03-w03-p03-010-bookkeeping}), \(\mathcal N\) is the union of
the elementary chain \(\langle\mathcal N_\alpha:\alpha<\lambda\rangle\).
By \((C_\lambda)\), \(\mathcal N_\alpha\preccurlyeq\mathcal N\) for every
\(\alpha\leq\lambda\); in particular \(\mathcal M\preccurlyeq\mathcal N\).

\emph{Saturation.}  Let \(A\subseteq\dom\mathcal N\) with
\(\abs A\leq\kappa\), and let \(p\) be a complete type over \(A\) in
finitely many variables consistent with \(\mathcal N\).  The domains
\(\dom\mathcal N_\alpha\), \(\alpha<\lambda\), increase and have union
\(\dom\mathcal N\).  By \cref{thm:p03-w03-p03-010-bookkeeping}, there is
\(\gamma<\lambda\) with \(A\subseteq\dom\mathcal N_\gamma\).  Since
\(\mathcal N_\gamma\preccurlyeq\mathcal N\),
\cref{thm:p03-w03-p03-010-type-facts}~(3) shows that \(p\) is a complete
type over \(A\) consistent with \(\mathcal N_\gamma\).  As
\(\mathcal N_{\gamma+1}\in W(\mathcal N_\gamma)\), some tuple \(\bar b\)
of \(\mathcal N_{\gamma+1}\) realizes \(p\) in \(\mathcal N_{\gamma+1}\).
As \(\gamma+1<\lambda\) and \(\mathcal N_{\gamma+1}\preccurlyeq\mathcal
N\), part~(3) of the same lemma shows that \(\bar b\) realizes \(p\) in
\(\mathcal N\).

Every set of cardinality below \(\kappa^+\) has cardinality at most
\(\kappa\), so \(\mathcal N\) is \(\kappa^+\)-saturated, and therefore
\(\kappa\)-saturated (\cref{def:p03-w03-p03-010-saturation-scope}).  The
statement of \cref{thm:saturated-extension} is the last conclusion.
\end{proof}
\end{annexitem}

\begin{annexitem}{annex:p03-w03-p03-010-size}{thm:p03-w03-p03-010-size}{Size of saturated structures and the finite case}
\begin{proof}
(1) Let \(\mathcal N\) be infinite and \(\kappa\)-saturated, and suppose
\(\abs{\dom\mathcal N}<\kappa\).  Put \(A=\dom\mathcal N\) and
\(p(x)=\{\neg(x=c_a):a\in A\}\).  A finite subset of \(p\) excludes
finitely many elements, and the infinite domain contains another one, which
realizes that subset.  By \cref{thm:p03-w03-p03-010-type-facts}~(1) and~(2),
\(p\) is contained in a complete type \(q\) over \(A\) consistent with
\(\mathcal N\).  Since \(\abs A<\kappa\), saturation gives \(b\in A\)
realizing \(q\); but \(\neg(x=c_b)\in q\) fails at \(b\).  So
\(\abs{\dom\mathcal N}\geq\kappa\).

(2) Let \(\mathcal M\) have exactly \(k\geq1\) elements, and let
\(\varepsilon_k\) be the \(\Lang\)-sentence
\[
  \exists x_1\cdots\exists x_k\Bigl(\bigwedge_{1\leq i<j\leq k}\neg(x_i=x_j)
  \land\forall y\bigvee_{i=1}^{k}y=x_i\Bigr),
\]
which uses only equality and is true in \(\mathcal M\).  If
\(\mathcal M\preccurlyeq\mathcal N'\), then \(\mathcal N'\entails
\varepsilon_k\), so \(\dom\mathcal N'\) has exactly \(k\) elements and
contains the \(k\) elements of \(\mathcal M\).  The domains are therefore
equal, and a substructure with the whole domain has the same
interpretations, so \(\mathcal N'=\mathcal M\).  By
\cref{thm:p03-w03-p03-010-saturated-extension}, \(\mathcal M\) has a
\(\kappa\)-saturated elementary extension, which is \(\mathcal M\) itself.
\end{proof}
\end{annexitem}

\begin{annexitem}{annex:p03-w03-p03-010-omega-example}{thm:p03-w03-p03-010-omega-example}{An \(\aleph_1\)-saturated elementary extension of \((\N,<)\)}
\begin{proof}
By \cref{thm:p03-w03-p03-010-saturated-extension} with \(\kappa=\aleph_0\),
\(\mathcal N\) realizes every complete type in finitely many variables over
every countable parameter set consistent with it, and
\(\N\subseteq\dom\mathcal N\).  The following \(\Lang(\N)\)-sentences are
true in \((\N,<)_\N\), for every \(n\in\N\):
\begin{align*}
  \sigma_n&:\ \forall x\bigl(\neg(c_n<x)\to(x=c_0\lor\cdots\lor x=c_n)\bigr),\\
  \rho_n&:\ \forall x\bigl(c_n<x\to(c_{n+1}<x\lor x=c_{n+1})\bigr),\\
  \tau&:\ \forall x\bigl(c_0<x\to\exists y\,(y<x\land
        \forall z\,(z<x\to(z<y\lor z=y)))\bigr),
\end{align*}
together with the axioms of a strict linear order and the atomic sentences
\(c_n<c_{n+1}\).  By \cref{thm:p03-w03-p03-010-type-facts}~(3) they are
true in \(\mathcal N_\N\).

(1) Let \(e\) be nonstandard and \(n\in\N\).  If \(\neg(n<e)\), then
\(\sigma_n\) puts \(e\) among \(0,\ldots,n\), contrary to assumption; so
\(n<e\).  In particular \(0<e\), and \(\tau\) gives \(y<e\) with every
\(z<e\) satisfying \(z\leq y\).  Nothing lies strictly between \(y\) and
\(e\), for such a \(z\) would satisfy both \(y<z\) and \(z\leq y\); so
\(y\) is an immediate predecessor of \(e\).  Suppose \(y=k\in\N\).  Since
\(k<e\), \(\rho_k\) gives \(k+1<e\) or \(e=k+1\); the latter is excluded
because \(e\) is nonstandard.  Then \(k+1<e\) gives \(k+1\leq y=k\), which
contradicts \(k<k+1\) in the strict linear order.  So \(y\) is
nonstandard.

(2) and (3).  Let \(E\) be a countable, possibly empty, set of nonstandard
elements, put \(A=\N\cup E\), which is countable, and let
\[
  p(x)=\{c_n<x:n\in\N\}\cup\{x<c_e:e\in E\}.
\]
A finite subset of \(p\) involves numbers \(n\leq m\) and finitely many
\(e_1,\ldots,e_k\in E\).  If \(k=0\), the element \(m+1\in\N\) realizes it,
because \(n<m+1\) is true in \((\N,<)\) and remains true in \(\mathcal N\).
If \(k\geq1\), let \(e\) be the least of \(e_1,\ldots,e_k\) in the linear
order, and let \(y\) be its immediate predecessor, nonstandard by (1).  By
(1), \(n<y\) for every \(n\leq m\), and \(y<e\leq e_i\) for each \(i\).  So
every finite subset of \(p\) is realized in \(\mathcal N\).  By
\cref{thm:p03-w03-p03-010-type-facts}~(1) and~(2), \(p\) is contained in a
complete type \(q\) over \(A\) consistent with \(\mathcal N\), and since
\(\abs A\leq\aleph_0\), some \(d\in\dom\mathcal N\) realizes \(q\).  Then
\(n<d\) for every \(n\in\N\), so \(d\notin\N\) by irreflexivity, and
\(d<e\) for every \(e\in E\).  With \(E=\varnothing\) this proves (2).  For
(3), a set \(E\) is coinitial among the nonstandard elements when every
nonstandard element is at least some member of \(E\); the element \(d\) is
below every member of \(E\) and so, by linearity, at least none of them.

(4) The structure \(\mathcal N\) is infinite, as it contains \(\N\), and it
is \(\aleph_1\)-saturated.  By \cref{thm:p03-w03-p03-010-size}~(1) with
\(\kappa=\aleph_1\), \(\abs{\dom\mathcal N}\geq\aleph_1\).
\end{proof}
\end{annexitem}

\begin{annexitem}{annex:p03-w03-p03-010-long-tuples}{thm:p03-w03-p03-010-long-tuples}{Realizing a type one coordinate at a time}
\begin{proof}
We construct \((b_\alpha)_{\alpha<\mu}\) in \(\mathcal N\) by
transfinite recursion so that, at stage \(\alpha\), the chosen tuple
realizes every member of \(p\) whose free variables have indices below
\(\alpha\). At stage zero these are sentences over \(A\), which hold
in \(\mathcal N_A\) because \(p\) is consistent with its complete
theory. At a limit stage take the union of the previous choices. Each
formula has finitely many free variables, so its indices are bounded
below that limit and its truth has already been ensured.

At a successor step put
\(A_\alpha=A\cup\{b_\beta:\beta<\alpha\}\). Since
\(\alpha<\mu\leq\kappa\), this parameter set has size less than
\(\kappa\), also when \(\kappa\) is singular: the union of two sets
of cardinality less than an infinite cardinal has cardinality less than
that cardinal. Consider all formulas in one variable \(z\) obtained
from members of \(p\) using variables among
\(\{x_\beta:\beta\leq\alpha\}\), by replacing \(x_\alpha\) by
\(z\) and earlier variables by their chosen values.

Every finite subset of this set is realized in \(\mathcal N\).
Indeed, take the corresponding finite conjunction \(\theta\) of
members of \(p\). The formula \(\exists x_\alpha\,\theta\) involves
only earlier variables. It belongs to \(p\): it is a logical consequence
of finitely many members of \(p\), and its negation would contradict
consistency, while completeness decides it. The recursion invariant
therefore makes that existential formula true of the earlier choices in
\(\mathcal N\). Its witness realizes the given finite subset.
By \cref{thm:p03-w03-p03-010-type-facts}, this one-variable set extends
to a complete type over \(A_\alpha\) consistent with \(\mathcal N\).
Saturation supplies \(b_\alpha\) realizing it. The invariant now holds
for all formulas using variables through index \(\alpha\).
Choice permits each selection, for example by fixing a well-order of
the domain and taking its first suitable witness. At the end every
formula of \(p\) is covered, proving realization of the whole type.
The empty variable family needs only the stage-zero observation.

The converse for the equivalent definition is immediate, because every
finite variable family has size less than the infinite \(\kappa\).
For the boundary example let \(M\) be a pure equality set of infinite
cardinality \(\kappa\). A complete finite-tuple type over
\(A\subseteq M\), \(|A|<\kappa\), specifies equality to elements of
\(A\), equalities among its coordinates, and which remaining classes
are distinct from all parameters. These exhaust the formulas: in a
quantifier-free formula write the Boolean combination of equalities in
disjunctive normal form. To eliminate an existential variable from a
conjunct, an equality to another variable or parameter permits
substitution. With no such equality, the only restrictions on that
variable are finitely many inequalities; an infinite set always has
a fresh value satisfying them. Contradictory conjuncts are discarded.
Induction on formulas therefore gives quantifier elimination for the
theory of an infinite equality set.
There are only finitely many fresh classes to realize, and
\(M\setminus A\) supplies distinct values for them. Thus \(M\) is
\(\kappa\)-saturated. The set
\(\{x_\alpha\neq x_\beta:\alpha<\beta<\kappa^+\}\) is finitely
satisfiable in \(M\), hence consistent with its theory by compactness,
and extends to a complete type in these variables. A realization would
inject \(\kappa^+\) into \(M\), contradicting \(|M|=\kappa\).
\end{proof}
\end{annexitem}
